\documentclass[11pt,a4paper]{article}
\usepackage[T1]{fontenc}
\usepackage[utf8]{inputenc}
\usepackage{lmodern}
\usepackage[margin=27mm]{geometry}
\usepackage{amsmath,amssymb}
\usepackage{booktabs,tabularx,longtable,array}
\usepackage{microtype}
\usepackage[hyphens]{url}
\usepackage[hidelinks]{hyperref}
\usepackage{enumitem}
\setlist{itemsep=3pt,topsep=5pt}
\setlength{\emergencystretch}{2em}
\setlength{\parskip}{4pt}
\setlength{\parindent}{0pt}
\newcommand{\ket}[1]{\lvert #1\rangle}
\newcommand{\bra}[1]{\langle #1\rvert}
\newcommand{\Tr}{\operatorname{Tr}}
\newcommand{\LocSig}{\mathsf{LocSig}}
\setcounter{tocdepth}{1}
\title{Locality, Coherent Alternatives, and Physical Records:\\
An Interpretive Account of Quantum Experiments in Modal Triplet Theory}
\author{Peter Nero}
\date{Version 1.0, 6 September 2026\\Unpublished interpretive manuscript}
\begin{document}
\maketitle
\begin{abstract}
Quantum experiments are often explained with words that change meaning during the explanation: locality becomes Bell factorization, a coherent alternative becomes another actual object, and a state update becomes a physical influence. This paper develops a qualified vocabulary and an interpretation-first account of the relevant experiments within Modal Triplet Theory. Beginning with individual photons and electrons, it distinguishes path coherence, localized detection, detector entanglement, one-arm monitoring, null records, many-particle effects, delayed conditioning, temporal correlations, and Bell nonfactorizability. It then connects these distinctions to a local-first constraint interpretation: actual events and records are generated by physical processes, while coherent source relations constrain their possible continuations. No fundamental universal wavefunction, additional spacelike source shortcut, or privileged role for consciousness is assumed. The account incorporates existing scoped MTT results on the canonical record measure, free local field algebra, compatible local descent, compression leakage, and the non-entailment of an ontology from operational data. Wigner's friend is analyzed through separate record-reading and coherent-reversal protocols, with actuality without automatic dephasing identified as an interpretive commitment rather than a derived universal law. The contribution is a unified conceptual account with reproducible finite examples, not a new derivation of every quantum experiment or a claim that the remaining physical apparatus and actualization problems are solved.
\end{abstract}

\section*{Version 1.0 Revision Note}
\begin{description}[style=nextline,itemsep=1pt,topsep=2pt]
\item[Supersedes:] No published paper. This manuscript consolidates the locality and quantum-ontology working notes.
\item[Reason:] The experimental analysis and local-first interpretation need one vocabulary.
\item[Resolution:] Definitions, experiments, MTT results, and interpretive commitments are distinguished.
\item[Retained content:] The earlier constraint-realism essay remains a separate companion.
\item[Remaining boundary:] Selected apparatus maps beyond the canonical domain and a universal objective-actualization law are not supplied here.
\end{description}

\newpage
\tableofcontents
\newpage
\section{The question is what the experiment warrants}
\label{sec:orientation}
A single detection event is easy to describe: a detector changes, a record is produced, and a particular location or channel is registered. The difficulty begins when that event is used to tell a story about everything that happened before it. A dot on a screen becomes evidence for a tiny classical object following one unobserved trajectory. An interference pattern becomes evidence for a second actual world. Correlation becomes an instantaneous message. None of these further identifications follows merely from the observed dot or pattern.

The purpose of this paper is not to replace accurate quantum calculations with a more appealing vocabulary. It is to keep four questions distinct. What was prepared and observed? What mathematical model predicts it? What ontology might that model support? Which additional MTT source construction, if any, selects the model? A standard quantum calculation can be correct before its MTT interpretation is derived. An interpretation can be intelligible without being uniquely forced by that calculation.

The proposed reading is a form of \emph{local-event constraint realism}. Physical events and records occur through interactions. Their possible continuations are constrained by phase, symmetry, admissibility, and composition structures that need not themselves be additional places. Quantum states represent these structures on a declared domain. One actual event need not mean that every coherent alternative was another actual object; conversely, calling alternatives modal cannot excuse losing their measurable phase relations.

The earlier essay on causal-base constraint realism develops the broader relation among closure repair, geometry, effective fields, and structural realism \cite{NeroConstraint}. The present paper has a different organizing question: how should that interpretation describe particular experiments without changing definitions halfway through? It is an experimental and conceptual companion, not a duplicate theorem source.

The reading order moves from individual excitations to correlations between systems, then to time and Wigner's friend. Modest equations identify distinctions that prose alone easily obscures. The appendix lists executable examples. The paper contains no new theorem/proof environments and does not take ownership of the imported fixed-point, recorder, or descent theorems.

\paragraph{Three kinds of statement.}
An \emph{established operational statement} belongs to standard physics or follows from an explicitly given ideal model. A \emph{scoped MTT result} is an accepted mathematical or selected-domain result with its stated hypotheses. An \emph{interpretive proposal} says what that structure might mean physically. These categories are explained in the surrounding prose rather than attached as administrative labels to every sentence.

\section{One photon and one electron}
\label{sec:species}
Photons and electrons share quantum features, but they are not interchangeable examples of the same classical substance. An electron is a charged, massive spin-one-half excitation. A photon is an electrically neutral, massless excitation of the electromagnetic field, with two physical helicities in the free theory. Polarization and electron spin can each provide two-state systems, but the similarity of their small state spaces does not equate their transformation laws, interactions, or statistics.

\begin{center}
\begin{tabular}{@{}>{\raggedright\arraybackslash}p{0.28\textwidth}>{\raggedright\arraybackslash}p{0.32\textwidth}>{\raggedright\arraybackslash}p{0.32\textwidth}@{}}
\toprule
Feature & Electron & Photon\\
\midrule
Free-particle character & Massive fermionic excitation & Massless bosonic excitation\\
Electric charge & Negative elementary charge & Zero\\
Spin information & Spin-one-half representation & Helicity \(+1\) or \(-1\)\\
Typical detector coupling & Charge deposition, scattering, ionization & Absorption or field-mediated detector excitation\\
Coherent alternatives & Spatial, spin, energy, or temporal modes & Spatial, polarization, frequency, or temporal modes\\
\bottomrule
\end{tabular}
\end{center}

For a nonrelativistic electron in a suitable regime, a spatial wavefunction is a useful approximation. For light, the general account uses quantum fields and a photodetection model. A scalar two-path photon amplitude is a controlled mode description, not a universal photon position wavefunction. An ideal one-photon state can have a vanishing mean electric field while retaining nonzero detection probabilities and interference. The relevant field correlations are not exhausted by the mean field \cite{Glauber}.

A detector click is a localized interaction. It does not establish that the excitation was a classical point object at every earlier instant. Likewise, extended coherence does not establish that energy was deposited continuously along every alternative. These are separate assertions about preparation, propagation, and interaction.

The MTT interpretation treats a particle as a persistent, physically realized response structure rather than a classical bead supplemented by a mysterious wave. Proto-spinor and shared-carrier encodings motivate a common language for some transformation and closure data, but they do not make photons into electrons or derive their masses from a two-state example \cite{NeroProto,NeroFoundation}. The species-specific field and detector maps remain essential.

No rest-frame story should be assigned to a photon. Null proper time does not create an observer perspective in which the photon sees its whole path at once. This tempting image does not supply an explanation of interference or entanglement.

\section{What exactly is meant by locality?}
\label{sec:locality}
The sentence "the theory is local" hides several logically different tests. We use a locality signature
\begin{equation}
\LocSig=(L_{\rm src},L_4^{\rm dyn},L_4^{\rm alg},
              B_{\rm fac},N_{\rm sig},M_{\rm ind}).
\end{equation}
Its entries refer to source-incidence locality, effective causal propagation, effective microcausality, Bell factorization, operational no-signalling, and setting independence. A signature must name its model, state class, operations, and domain. Its entries may be proved, conditional, violated, or undecided; an aspirational string of ones and zeros is not an experimental result.

\subsection{Differential support and causal propagation}
A local differential equation uses a field and finitely many derivatives at the same point. This is a useful structural property but not a sufficient causality theorem. Hyperbolicity, constraints, domains, and the propagation problem also matter. An ordinary heat equation is differential yet has noncompact spatial response at every positive time.

For a relativistic retarded response written as \(G_{\rm ret}(x,y)\), the relevant support condition is
\begin{equation}
G_{\rm ret}(x,y)=0\quad\hbox{unless }y\in J^-(x).
\end{equation}
Here \(J^-(x)\) is the causal past of the response point. Integral notation alone does not make an equation acausal: a kernel can have strictly causal support. Conversely, a spatially extended equal-time expression requires investigation rather than automatic identification with a controllable signal.

\subsection{Algebras and experimental marginals}
Microcausality concerns observables in spacelike-separated regions:
\begin{equation}
[A,B]=0,\qquad A\in\mathcal A(O_A),\ B\in\mathcal A(O_B),
\quad O_A\perp O_B.
\end{equation}
Odd fermionic fields use graded commutation; even physical observables commute ordinarily. Commutation does not make the state a product. A state on commuting subsystems can still be entangled.

No-signalling instead concerns an intervention and a probability:
\begin{equation}
\sum_b p(a,b\mid x,y)=p(a\mid x).
\label{eq:nosignal}
\end{equation}
The remote setting \(y\) does not alter the local marginal. To infer this from an algebraic model, the allowed operations must have the required localization and normalization. An arbitrary projector written on a large Hilbert space is not automatically an admissible relativistic measurement \cite{FewsterVerch}.

\subsection{Source-incidence locality}
Let \(N_S^-(e)\) be a declared primitive predecessor neighborhood of a source event or configuration \(e\). An illustrative deterministic law is source-incidence local when
\begin{equation}
s'(e)=F_e\!\left(s|_{N_S^-(e)}\right).
\end{equation}
This defines dependence relative to the source relation. It does not identify that relation with the effective light cone. A restriction-compatible local map is another useful mathematical notion, but a local infinitesimal rule need not remain neighborhood-supported after finite evolution.

No additional source shortcut between spacelike-separated effective locations is assumed in this paper. Entangled states and local detector interactions already provide the operational architecture that must be explained. If a later source completion introduces a different causal structure, its physical consequences must be stated independently.

The phrase constraint direction is deliberately neutral about compactness and physical spatial extent. Some MTT data include compact phase circles. Other internal variables may be noncompact. Neither property alone makes a variable another ordinary spatial direction.

\section{Two slits: coherence before a record}
\label{sec:slits}
The central experiment does not require two particles to meet. Electrons prepared and detected individually can build a double-slit distribution from successive localized events \cite{Bach}. Single-photon experiments likewise distinguish interference from the anticorrelation characteristic of the relevant single-photon preparation \cite{Grangier}. The preparation and detection qualifications matter: a mean optical fringe alone also occurs with classical coherent light.

Write the two coherent channel contributions at screen coordinate \(q\) as
\begin{equation}
A(q)=\alpha\psi_L(q)+\beta e^{i\phi}\psi_R(q).
\end{equation}
The detection density is proportional to
\begin{align}
|A(q)|^2={}&|\alpha|^2|\psi_L(q)|^2+
|\beta|^2|\psi_R(q)|^2 \nonumber\\
&+2\operatorname{Re}\!\left[
\alpha^*\beta e^{i\phi}\psi_L(q)^*\psi_R(q)\right].
\label{eq:interference}
\end{align}
The last term is the observable consequence of the phase relation. It is not an additional stream of particles.

An intuitive picture is to add two arrows, including their directions, before squaring the resulting length. Sometimes they reinforce, sometimes they cancel. Adding the two squared lengths instead gives a different law. The arrows are a picture of amplitudes, not little arrows physically attached to classical objects.

If the input is an incoherent path mixture, its probability has only the first two terms. Therefore "we have not selected a path yet" is not enough to explain the experiment. The alternatives must remain coherently related. In this paper, \emph{coherent modal plurality} names that interpretive reading of superposition.

There is no inference here that the electron was two actual electrons or that two actual universes were required. There is also no inference that it secretly chose one classical path while all other physically effective structure disappeared. A definite-trajectory completion would need some additional structure retaining the phase dependence of both alternatives.

For a balanced two-mode interferometer, the output probabilities are
\begin{equation}
p_\pm=\frac{1\pm\cos\phi}{2}.
\end{equation}
The same excitation can therefore produce different localized output statistics when the relative phase changes. Wave and particle language describe different aspects of this one experiment: coherent propagation and localized exchange.

\section{A detector at only one slit}
\label{sec:onearm}
This case exposes a particularly persistent ambiguity. A device at the left slit can change the final interference pattern without changing the probability of a right-arm population measurement performed before recombination. The two observations concern different quantities.

Let a nondestructive marker correlate the alternatives with apparatus states:
\begin{equation}
\ket{\Psi}=\alpha\ket{L}\ket{D_L}
            +\beta e^{i\phi}\ket{R}\ket{D_R}.
\end{equation}
For a screen readout that does not retain the marker, define
\(\gamma=\langle D_L|D_R\rangle\). The cross term in
Eq.~\eqref{eq:interference} is multiplied by \(\gamma\); the path populations remain \(|\alpha|^2\) and \(|\beta|^2\). Orthogonal marker states suppress that cross term in the unconditioned screen distribution.

Even when only the left arm interacts actively with a device, the apparatus states "interacted" and "did not interact" can distinguish the alternatives. Nothing in this calculation requires a message to cross from the left slit to a separately prepared right-path particle.

In the pure, balanced marker model the visibility is \(|\gamma|\), and optimal which-way distinguishability is \(\sqrt{1-|\gamma|^2}\). More generally their complementarity is an inequality, not a claim that a mechanical kick is always necessary \cite{Englert}.

\subsection{What happens to a known right-path input?}
For a concrete imperfect monitor, use
\begin{equation}
M_0=\sqrt{1-\eta}\ket L\bra L+\ket R\bra R,\qquad
M_1=\sqrt{\eta}\ket L\bra L,\quad 0\leq\eta\leq1.
\label{eq:monitor}
\end{equation}
The effects sum to the identity. The unread channel preserves the diagonal path probabilities and multiplies their off-diagonal coherence by \(\sqrt{1-\eta}\). A right-only input \(\ket R\bra R\) is unchanged.

Accordingly, two claims must not be substituted for one another:
\begin{quote}
Changing a one-arm coupling changes a coherently recombined distribution.

Measuring a left-path particle changes a different, independently prepared right-path particle.
\end{quote}
The first is realized by the model. The second is not. In the unmarked experiment, the screen events were not partitioned into experimentally certified prior left and right trajectories.

\subsection{A null record is an outcome of a protocol}
For balanced input, the no-indication probability in Eq.~\eqref{eq:monitor} is \(1-\eta/2\). The conditional right-path probability is \(1/(2-\eta)\). Detector silence identifies the right path with certainty only in the ideal \(\eta=1\) version with the other relevant assumptions satisfied. Inefficiency, losses, false indications, and the monitoring window can change the inference.

At \(\eta=9/25\), the unread output probabilities after balanced zero-phase recombination are \(9/10\) and \(1/10\). The no-indication subset occurs with probability \(41/50\), and within it the two outputs have probabilities \(81/82\) and \(1/82\). These exact fractions are not fitted physical constants; they illustrate a chosen apparatus efficiency.

An absorber with the same no-absorption operator is a different instrument. It removes the absorbed branch from the screen ensemble. In the same example the screen probabilities are \(81/100\) and \(1/100\), with absorption probability \(18/100\). Renormalizing only the survivors must not conceal the missing particles.

Interaction-free measurement uses interferometric alternatives and a suitably informative null or dark-port outcome. Its force is not that no physical experimental arrangement was present, nor that every run avoids interaction. The proposed absorber changes the available process, and the outcome statistics distinguish it \cite{ElitzurVaidman}.

\subsection{The MTT reading}
The source-supported relation between alternatives has changed because an apparatus has become correlated with them. The reduced screen description no longer retains the same coherence. This is a statement about a joint relation and a specified readout, not about a photon discovering that a person is watching.

The interpretation directs attention to continuation conditions rather than a choice between classical wave and classical particle substances. It is not uniquely quantum-mechanical evidence for MTT: standard instruments already calculate the effect. MTT's additional task is to obtain those instruments from its selected physical source.

\section{One run, many runs, and many particles}
\label{sec:number}
In one ideal single-excitation run, one registered detection need not reveal a fringe. A fringe is a distribution estimated across runs. For independent preparations with probability \(p\) of a specified output, the count over \(N\) runs has mean \(Np\) and variance \(Np(1-p)\). The successive particles need not communicate to draw that distribution.

Independence is an assumption about preparation and apparatus memory, not a consequence of low count rate alone. Drift, a retained environment, or intentionally correlated preparations can connect trials. These effects require an enlarged process model, not the declaration that all particles always influence one another.

Many-particle experiments introduce additional questions. Identical-particle exchange, entanglement between accessible subsystems, and ordinary classical correlations are different structures. Two indistinguishable photons entering opposite inputs of a balanced beam splitter exhibit the Hong--Ou--Mandel effect under the appropriate mode overlap \cite{HongOuMandel}. Fermionic exchange gives different behavior, but saying that all electron pairs simply repel quantum-mechanically would ignore spin, state preparation, and interactions.

First-order fringes alone do not establish single-photon preparation. For an ideal single-mode number state with \(n>0\),
\begin{equation}
g^{(2)}(0)=\frac{\langle n(n-1)\rangle}{\langle n\rangle^2}
         =1-\frac1n.
\end{equation}
The ideal one-photon value is zero, whereas a coherent state has value one. Actual measurements require detector and background analysis. Particle-number statistics and path coherence answer different questions \cite{Glauber,Grangier}.

The MTT interpretation must retain a distinction between an excitation, an experimental trial, a field mode, and a subsystem. Counting modes or source coordinates is not counting actual particles. Nor is a tensor factor automatically another universe.

\section{Entanglement is not just a shared past}
\label{sec:entanglement}
Two systems may share an origin and remain only classically correlated. A separable joint state has the form
\begin{equation}
\rho_{AB}=\sum_k p_k\,\rho_A^{(k)}\otimes\rho_B^{(k)}.
\end{equation}
The classical label \(k\) can encode a common preparation. Entanglement means that no such decomposition exists for the declared subsystems. It is a stronger property than common ancestry, and it is defined relative to a composition structure.

Bare single-particle path interference does not require entanglement between two separately prepared particles. Entanglement enters the marker example when nontrivial path alternatives become correlated with nonparallel pure marker states. There is also a mode-entanglement convention for a state such as \((\ket{1,0}+\ket{0,1})/\sqrt2\). Its subsystem choice should be named rather than used to turn a one-photon experiment into a two-photon experiment.

In MTT language, \emph{shared-source provenance} describes common ancestry; \emph{shared-source coherence} proposes a phase-sensitive mechanism; \emph{entanglement} describes the resulting nonseparability. These expressions may be connected by a physical model, but they are not synonyms.

A useful intuition is that individual descriptions do not exhaust the physically relevant relation. The two local reduced states of an entangled pair can be maximally mixed while their joint statistics are highly constrained. This is not a missing classical answer stored independently inside each particle.

Entanglement does not imply Bell violation for every measurement choice. Entangled mixed states can admit Bell-local descriptions for specified measurement classes \cite{Werner}. A claim about preparation entanglement, one about a Bell test, and one about signalling therefore need separate evidence.

\section{Bell: which classical explanation fails?}
\label{sec:bell}
Let \(x,y\) be settings, \(a,b\) outcomes, and \(\lambda\) the complete relevant source variable of a proposed classical completion. Bell factorization and setting independence are
\begin{align}
p(a,b\mid x,y,\lambda)&=p(a\mid x,\lambda)p(b\mid y,\lambda),\\
\rho(\lambda\mid x,y)&=\rho(\lambda).
\end{align}
They give a common probability mixture of local response models. For the usual dichotomic CHSH experiment, that class satisfies \(|S|\leq2\). The standard singlet model reaches \(2\sqrt2\) for appropriate settings while retaining equal local marginals. Bell's argument and modern experiments constrain the factorized independent-setting model, not a slogan that all correlations must be signals \cite{Bell,Hensen}.

An upper or source description is not exempt from this argument. If it supplies a complete classical \(\lambda\), independent settings, and factorized effective responses, changing the name of \(\lambda\) to a bundle configuration does not change the bound. A proposal must identify the mathematical property that fails.

In a candidate quantum-source architecture the effective observables can be microcausal, operational marginals can obey Eq.~\eqref{eq:nosignal}, and the correlations can lack a Bell-factorized independent-setting representation. Such an architecture is familiar at the operational quantum level. A selected MTT construction of all its states and instruments is an additional claim.

Deterministic source evolution does not automatically produce deterministic projected outcomes. If a completion does posit deterministic outcomes and setting independence, it must confront how its response law escapes Bell factorization. Source determinism alone does not settle this either way.

\subsection{Measurement dependence and superdeterminism}
Measurement dependence is the statistical statement
\begin{equation}
\rho(\lambda\mid x,y)\ne\rho(\lambda).
\end{equation}
It is not defined by a metaphysical opinion about free will. Common ancestry of two measured systems does not establish it. Common ancestry of the source and the setting generators may produce it, but the relevant joint distribution must actually be evaluated.

The term superdeterminism is commonly used for deterministic Bell-foundations accounts built around such dependence. A retrocausal or time-symmetric model can also be measurement-dependent without having the same causal story. These alternatives should not be silently merged.

The present interpretation does not add measurement dependence or a spacelike source shortcut as an axiom. It also does not announce a completed six-entry locality signature for physical MTT. Its baseline is to preserve independently chosen laboratory controls and no-signalling while investigating the actual source composition law \cite{NeroBell}.

\section{Time: intervention, memory, and delayed conditioning}
\label{sec:time}
The distinction between an intervention and a condition is as important along the time axis as it is between two laboratories. An earlier interaction can alter a later distribution. A later label can partition earlier records. These are not the same operation.

\subsection{Delayed choice and quantum erasure}
A delayed-choice interferometer changes the readout arrangement after an excitation has entered it. The appropriate model includes that time-dependent apparatus. It need not say that the choice travels back to make the excitation previously a wave or previously a particle \cite{Jacques}.

For an ideal path-marked state, a marker measurement in a complementary basis produces conditional path states with relative signs \(+\) and \(-\). Their screen distributions have complementary fringes, but the unconditioned sum does not. An eraser experiment can use a later marker label to identify these subsets; the earlier unsorted detections are not altered \cite{Kim}.

For \(\ket{\psi_\pm}=(\ket L\pm\ket R)/\sqrt2\), the relevant identity is
\begin{equation}
\tfrac12\ket{\psi_+}\bra{\psi_+}+
\tfrac12\ket{\psi_-}\bra{\psi_-}
=\tfrac12\bigl(\ket L\bra L+\ket R\bra R\bigr).
\end{equation}
Destroying a classical file is not the coherent marker measurement in this equation. Nor does seeing fringes in a selected subset imply that an observer could have sent a chosen message into the earlier unselected record list.

\subsection{Temporal coherence and quantum memory}
A single excitation can occupy coherent early and late temporal modes. Their relative phase becomes observable when the apparatus overlaps the alternatives appropriately. That is temporal-mode coherence, not a particle consulting a completed future.

Memory arises when the variables retained in a reduced description omit information relevant to later statistics. In quantum processes the omitted information can include system-environment correlations and the effects of interventions. An enlarged state description can sometimes restore a Markov property, but a proposed sufficient description must be demonstrated \cite{Pollock}.

For a finite causally ordered instrument model with freely specified external settings, completeness of every future instrument gives
\begin{equation}
\sum_{r_2,\ldots,r_n}p(r_1,\ldots,r_n\mid x_1,\ldots,x_n)
=p(r_1\mid x_1).
\end{equation}
This is an earlier-marginal statement under the model's causal and preparation assumptions. It allows ordinary forward disturbance. Conditioning on a later outcome is deliberately a different calculation.

\subsection{Temporal Bell tests and repeated monitoring}
For three dichotomic values with one common joint distribution,
\(Q_1Q_2+Q_2Q_3-Q_1Q_3\leq1\) pointwise. The Leggett--Garg use of this bound concerns macrorealism and noninvasive measurability, not the identical assumptions of a spatial Bell experiment \cite{LeggettGarg}.

In an ideal rotating-qubit example, separate pair-measurement protocols can give \(C_{12}=C_{23}=1/2\) and \(C_{13}=-1/2\), hence \(K=3/2\). If an intervening projective measurement is actually inserted, it changes the protocol and can change the last correlation. One must not pretend that the three pair experiments already supplied an undisturbed triple of values \cite{NeroTemporal}.

Repeated ideal monitoring can also alter survival. A standard Zeno example has survival \(\cos^{2N}(\theta/N)\), tending to one at fixed \(\theta\). The monitoring is a physical sequence, not human attention holding an object in place.

MTT's auxiliary closure-repair parameter is not automatically this physical clock. A heat semigroup used in selection or regularization must not be reinterpreted as instantaneous physical propagation without a separate bridge.

\section{Uncertainty, gauge comparison, and constraint directions}
\label{sec:gauge}
Uncertainty is not one single assertion about inaccurate instruments. For suitable states and self-adjoint operators on the necessary domains, the Robertson bound reads
\begin{equation}
\Delta A\,\Delta B\geq\tfrac12|\langle[A,B]\rangle|.
\end{equation}
It concerns state spreads. Measurement disturbance and simultaneous-measurement accuracy require their own definitions. Fourier-dual localization is an especially useful way to understand position-momentum tradeoffs, without imagining that a poorly manufactured detector is the whole explanation \cite{Robertson}.

MTT may read incompatibility as a limit on which effective sharp-property requirements are jointly admissible. Its compression results provide exact examples in which an effective commutator depends on excursions into an omitted sector. That is a mechanism worth interpreting, not proof that every physical uncertainty relation has already been sourced.

Gauge structure concerns comparison of local representatives. A change of gauge can alter a description without altering its physical content; a connection specifies comparison, and holonomy can carry observable information. The Aharonov--Bohm phase illustrates that local field-strength values along accessible paths need not exhaust gauge-invariant experimental data \cite{AharonovBohm}.

The MTT shared-circle proposal is naturally read as common phase or holonomy structure, not a wire carrying instantaneous commands between outcomes. Sharing a phase convention does not itself create entanglement. Nor does ordinary \(U(1)\) geometry automatically derive every non-Abelian gauge group. The source action, its stabilizer, the observable kernel, and a faithful quotient must be specified.

Constraint coordinates consequently require role tests. Are they gauge representatives, internal physical modes, deformation parameters, or actual additional causal directions? This paper does not decree that all extra dimensions in all theories are unreal. It proposes that coordinate counting alone cannot answer the physical question \cite{NeroConstraint,NeroFoundation}.

\section{From closure to operators: what the interpretation imports}
\label{sec:repair}
The phrase closure repair is meaningful only when a law or a sufficiently typed mathematical construction accompanies it. In one important structural setting a defect map \(\Phi\) and a chosen metric define a nonnegative cost \(\mathcal C=\|\Phi\|^2/2\). Negative-gradient flow reduces that cost, and its stationary configurations are candidates for stable closure.

At a zero-defect configuration, the repair Hessian has the normal-square form \(D\Phi^\dagger D\Phi\). Away from zero defect, additional terms can survive. Under the required operator hypotheses a positive self-adjoint tangent operator supports heat-semigroup and spectral-calculus constructions. These are imported structural results, not proofs that a universal physical Hamiltonian has been selected \cite{NeroRepair}.

A frequently useful retained operator is schematically
\begin{equation}
B_{\rm adm}=P\,\chi(A)e^{-\tau A}\chi(A)\,P.
\end{equation}
Here \(A\) is the appropriately realized tangent operator, \(\chi\) a specified spectral filter, and \(P\) the retained-sector projector. Its filtered spectral description and, when the required representation exists, its integral kernel describe one operator in different ways. The former suggests modes; the latter suggests a response between localized test data. This helps explain why wave-like and localized descriptions can coexist.

It does not yet derive a detector click, a Born law, or the physical meaning of \(\tau\). An arbitrary semigroup need not possess the regular integral kernel being imagined. A positive repair square is also not automatically the signed Lorentzian action governing physical propagation. The interpretation gains precision by preserving these boundaries rather than calling every operator a heat kernel.

Three reductions must remain distinct. A representation change gives another description of the same state. A restriction or compression discards accessible structure or changes products. A measurement instrument gives outcome-indexed state transformations and probabilities. The word projection cannot substitute for the differences among them.

\section{Results that give this account more than a vocabulary}
\label{sec:results}
The following results are incorporated from their owners, not reproved here. Their role is to identify where the interpretive picture already has mathematical content. The reproducibility section provides exact result identifiers and source locations.

\subsection{An exact canonical record law}
The current MTT canonical binary one-anchor recorder has a pointer/Fock dilation, a stopped completely positive instrument, and a classical first-record output measure on its commuting nondemolition output algebra. Its first-record probability and second-moment descent are exact on that declared domain, without adding a separate Born axiom, fitted output probability, or extra stochastic primitive there \cite{NeroQM,NeroResults}.

This is a real accomplishment. It does not mean that every laboratory apparatus has already been derived from the same action, nor that the existence of a probability law selects one ontically actual sample. A classical output algebra, a selected normal state, and a measure have definite mathematical roles; they should not be silently renamed a universal solution of the measurement problem.

\subsection{One record channel need not specify one instrument}
The imported recorder work also distinguishes instruments with different retained records despite the same unread channel. Discarding the output label can erase the distinction between an informative measurement and another realization of the same average transformation. Thus decoherence alone does not specify all the physical content of an experiment.

This is directly useful for one-arm monitoring and Wigner's friend. We must say which record is retained, which environment is accessible, and which operation is later performed. An averaged density matrix is not a complete experimental protocol \cite{NeroMeasurement,NeroResults}.

\subsection{Local field structure at a stated free-theory tier}
The selected framed q79 construction includes a free massless twisted Dirac CAR net with graded locality, an Einstein-local even observable subnet, causal propagation, time-slice structure, and a nonempty appropriate state space \cite{NeroResults}. It is therefore incorrect to describe all MTT locality as only a philosophical wish.

It would be equally incorrect to promote this free construction into a complete interacting theory of every electron, photon, and detector used above. The present paper preserves the free-net result while stating what the experimental extension asks for.

\subsection{Local-first descent and the meaning of global data}
For the imported finite local-system construction, local coherent representatives satisfy
\begin{equation}
(d_0\psi)_{ij}=\psi_j-U_{ij}\psi_i,\qquad \Gamma=\ker d_0.
\end{equation}
Compatible families are determined by overlap conditions and holonomy constraints, not by multiplying copies of one chart into separate universes.

The central-twist example is especially instructive. A cycle action \(\omega I\), with \(|\omega|=1\) and \(\omega\ne1\), has no nonzero fixed vector, while conjugation by it fixes all endomorphisms. The displayed vector equalizer has dimension zero and its endomorphism equalizer dimension nine. Rays and operator data can therefore remain meaningful when one chosen vector-level compatibility condition fails.

This does not prove that every global mathematical state is absent. It concerns a specified finite local-system equalizer, not arbitrary smooth sections of an arbitrary bundle. It also does not guarantee that all local marginals possess a common state extension \cite{NeroResults,NeroRepair}.

\subsection{Compression can create incompatibility}
For commuting upper self-adjoint operators \(\widetilde A,\widetilde B\) and an orthogonal projector \(P\), write
\(L_A=(I-P)\widetilde A P\) and \(L_B=(I-P)\widetilde B P\).
Under the required product and domain assumptions, the exact identity is
\begin{equation}
[P\widetilde A P,P\widetilde B P]=L_B^\dagger L_A-L_A^\dagger L_B.
\end{equation}
This locates an effective commutator in leakage through the excluded sector. It also prevents an invalid inference: upper commutation does not survive every arbitrary compression. Compatible reducing or product-preserving structures require separate verification \cite{NeroConstraint,NeroResults}.

\subsection{Operational data do not force many actual worlds}
At a declared canonical checkpoint, the imported ontology countermodel has record weights
\begin{equation}
(p_{\rm ready},p_P,p_Q)=\frac1{448}(1,149,298).
\label{eq:448}
\end{equation}
A uniform 448-atom sample space partitioned into sets of sizes \(1,149,298\) gives exactly those probabilities with one record label per atom. An exact unitary dilation gives the same operational weights. Distinct ontological completions therefore share the retained record data \cite{NeroResults}.

This establishes a scoped non-entailment result, not a physically selected hidden-variable theory of all experiments. The checkpoint does not assign simultaneous values to every incompatible projector. It proves neither that many-worlds interpretations are false nor that one actual history has been selected by MTT dynamics.

These distinctions are the paper's principal use of real MTT results. They support a disciplined single-world interpretation while preventing the interpretation from claiming more than its source.

\section{A proposed single-world ontology}
\label{sec:ontology}
The proposed picture has three levels. Actual events are particular physical occurrences. Coherent source relations constrain the responses that interactions can produce. State descriptions encode those relations for specified experimental domains. The levels are related, but none is a synonym for the others.

A coherent alternative is not automatically another actual person, apparatus, or universe. Equally, it cannot be reduced to ignorance about a classical fact if its relative phase affects a later experiment. Not actual does not mean mathematically removable without consequences.

In this interpretation, measurement is ordinary physical record formation. It is not an act by which consciousness confers existence. The special mathematical treatment arises because one asks about a recorded outcome and the state conditioned on it, rather than only about an averaged evolution.

The proposed account of actuality makes a further distinction: a definite local event need not be identical to irreversible physical dephasing of every coherent source relation. This is a substantive interpretive commitment. It must eventually receive a source law if it is to function as a completed physical explanation.

A simple counterexample explains why fixed-point uniqueness alone does not supply that law. A uniquely selected density operator or retained sector can still assign nontrivial probabilities to several detector outcomes. Uniqueness of the state solves a different problem from uniqueness of the outcome in a run. A basin may specify a stable response, but its existence does not automatically generate the correct ensemble weights.

The current corpus includes an explicitly conditional marked-Poisson construction in which supplied quadratic channel weights determine a first-record race. It is a constructive candidate for actualization, not an adopted axiom or a consequence of the shared circle alone. Its quadratic coupling and rate assumptions remain part of the hypothesis. This paper does not silently add that primitive to MTT \cite{NeroQM}.

\section{Wigner's friend, told from inside the world}
\label{sec:wigner}
The friend and Wigner are ordinary physical participants. Neither stands outside the universe. The interpretation assumes no independently existing universal wavefunction that one participant alone can read correctly.

Begin with a spin-like two-state system and a memory. An ideal reversible coupling produces
\begin{equation}
\ket{\Psi_+}_{SF}=
\frac{\ket{0}\ket{F_0}+\ket{1}\ket{F_1}}{\sqrt2}.
\label{eq:friend}
\end{equation}
The finite quantum-memory model is not itself a demonstration that a conscious observer can be coherently reversed.

\subsection{The friend records an outcome}
In the proposed one-world reading, a local interaction produces a particular record, say \(0\). This event does not await Wigner's attention. Information about it becomes available to Wigner through further physical coupling.

Wigner's coherent description, when the stated isolation and control assumptions hold, describes the joint response structure of the finite experiment. It need not describe two actual friends. Nor is it merely ignorance about a physically collapsed ensemble. Those two interpretations of the density data lead to different possible continuations.

The friend's statement "my record is 0" and Wigner's statement "this preparation supports a coherent joint test" are not automatically contradictory statements about the same mathematical object. The interpretation must explain how actual records and coherent response data coexist; it must not solve the problem merely by changing what the word state means.

\subsection{Wigner reads the record}
If Wigner performs an appropriate repeatable pointer readout while the record is retained, his record agrees with it. This is another physical interaction extending the record. There is no need for realities to merge or for a global update to have travelled instantaneously.

A conditional state based on a known record is not the same as an unconditioned state. Nevertheless, saying only that the observers have different information is insufficient for the full thought experiment, because Wigner can be imagined performing a different operation.

\subsection{Wigner attempts coherent reversal}
In a two-qubit realization, a controlled-NOT correlates the system and memory. Applying it again undoes the coupling:
\begin{equation}
\frac{\ket{00}+\ket{11}}{\sqrt2}\longmapsto\ket{+}\ket0.
\end{equation}
An \(X\) measurement of the system then gives \(+\) with probability one. If the physical state had instead become the dephased ensemble
\begin{equation}
\rho_{\rm mix}=\tfrac12\ket{00}\bra{00}+
\tfrac12\ket{11}\bra{11},
\end{equation}
the same sequence gives probability one half.

In the proposed actuality-without-automatic-dephasing reading, coherent source alternatives can matter to this later response without having been other actual worlds. The earlier event is not declared never to have happened; the distinguishable pointer record is no longer retained after the ideal reversal. This is a possible interpretive account, not an already derived universal event dynamics.

A physically received copy of the record changes the experiment. If an uncontrolled environment retains orthogonal states \(E_0,E_1\), the restricted system-memory description has lost the relevant coherence. More generally, in the balanced zero-phase example, a joint coherence projector has probability
\begin{equation}
p_+=\frac{1+\operatorname{Re}\langle E_0|E_1\rangle}{2}.
\end{equation}
Wigner cannot recover the original interference merely by forgetting what he learned or deleting a classical file. Reversing all relevant copies, where physically possible, is a different control operation. Preparing fresh coherence after a reset is also not recovery of the original recorded alternatives.

\subsection{No global state: what is and is not rejected}
A fundamental state of the entire universe is different from a compatible quantum description of one finite laboratory. The latter can be built from local preparations and causal interactions. The local-first interpretation rejects the necessity of the former; it does not automatically prohibit the latter.

Absence of a global vector representative is also different from absence of rays, operator data, or a state on an observable algebra. The MTT descent example in Section~\ref{sec:results} already demonstrates why the distinction matters.

Finally, a joint quantum state is not a common classical table of definite values for every incompatible measurement. Contextuality can obstruct the latter without preventing the former. A sheaf over measurement contexts and a sheaf over source charts have different bases and cannot be identified by vocabulary alone \cite{AbramskyBrandenburger}.

The strongest ordinary-language summary is therefore:
\begin{quote}
An event can be actual locally without being information already possessed everywhere. Coherent alternatives can constrain later responses without being other actual worlds. Wigner's readout and his attempted reversal are distinct physical continuations.
\end{quote}

\subsection{The extended experiments still constrain the account}
Frauchiger--Renner and Local Friendliness arguments investigate stronger combinations of observer reasoning, outcomes, and quantum control than ordinary record comparison. Their assumptions must be stated rather than reduced to the slogan that quantum mechanics contradicts itself. The Local Friendliness result constrains combinations of absolute observed events, locality, and setting independence when the relevant observer-scale quantum control and correlations are admitted \cite{FrauchigerRenner,Bong}.

Denying a fundamental universal state is not by itself an identified escape from that finite-experiment result. A completed MTT interpretation must specify which premise fails, which operation is unavailable, or which prediction differs. This paper offers a coherent vocabulary and a candidate ordinary-friend reading; it does not claim that erasing a record solves all extended no-go theorems.

\section{How the interpretation helps, and where it does not}
\label{sec:comparison}
The conceptual gain is to stop demanding that a quantum object be either a classical little body or a classical distributed substance. A physical excitation has a response structure. Its coherent alternatives determine possible interference; its coupling to a detector produces localized exchange and records.

Likewise, a joint relation need not be an exchanged message. A later condition need not be an earlier cause. A fiber coordinate need not be another location. These distinctions can make quantum experiments more intuitive, but they are also available in careful standard quantum accounts. MTT should not claim ownership of clarity already present in quantum instruments, algebraic locality, or decoherence theory.

What could be distinctive is upstream unification. The same independently selected source structure might explain why the appropriate phase relations, field algebras, admissible states, detector laws, and gauge comparisons fit together. The scoped results in Section~\ref{sec:results} provide parts of that architecture. They do not yet prove that the entire experimental family has one selected physical realization.

This view is close in spirit to process and structural forms of realism: events and lawful relations matter more than an inventory of intrinsic classical properties. It is not automatically equivalent to relational quantum mechanics, a modal interpretation, a pilot-wave theory, objective collapse, or Everettian branching. It does not dismiss those approaches by renaming their unresolved questions.

One actual world does not require a completed classical answer sheet for every unperformed experiment. Nor does the rejection of such an answer sheet imply that records depend on human opinion. The task is to model actual interactions and their compatibility without importing unjustified counterfactual identifications.

\section{A working glossary and recurring misconceptions}
\label{sec:glossary}
\begin{longtable}{@{}p{0.27\textwidth}p{0.67\textwidth}@{}}
\toprule
Expression & Meaning used here\\
\midrule
\endfirsthead
\toprule
Expression & Meaning used here\\
\midrule
\endhead
Coherent modal plurality & Alternatives whose relative phases remain effective in the relevant experiment; not a classical mixture.\\
Entanglement & Nonseparability for a declared subsystem composition; not merely a common origin.\\
Shared-source provenance & Common source ancestry, without an automatic claim about coherence or setting dependence.\\
Physical record & An accessible correlation produced by interaction; awareness is not a defining ingredient.\\
Actual event & A particular occurrence in the proposed ontology, distinguished from its probability law.\\
Source-incidence locality & Dependence on a specified primitive source neighborhood.\\
Effective causal locality & Physical response constrained by the effective causal cone.\\
Microcausality & Commutation of spacelike physical observable algebras.\\
Bell nonfactorizability & Failure of the declared Bell-factorized independent-setting representation.\\
No-signalling & Invariance of local marginals under allowed remote setting changes.\\
Measurement dependence & Statistical dependence of the relevant source variable on Bell settings.\\
Constraint direction & A source degree of freedom whose physical role requires specification; not automatically another place.\\
Local-first state & Local or compatible finite-domain response data, without a postulate of a fundamental universal state.\\
\bottomrule
\end{longtable}

\paragraph{Several corrections should travel together.}
A click does not certify an unmeasured trajectory. A lost fringe does not prove a mechanical kick. A detector at one arm does not thereby disturb independently prepared particles in the other. No click does not always mean the other path. Absorption and unread nondestructive marking are not the same channel on surviving particles.

Many trials do not require particle-particle communication. Many particles do not automatically imply entanglement. A single-photon state is not a weak classical wave merely because both can show mean fringes. A later eraser label does not redraw earlier detections. Temporal Bell and spatial Bell experiments do not use identical assumptions.

An upper shared circle does not automatically solve Bell's factorization problem. A unique fixed point does not automatically select one detector outcome. A heat-kernel parameter is not automatically physical time. Neither a global Hilbert notation nor the number of nonzero coefficients counts actual universes.

\section{The physical comparison that would decide more}
\label{sec:frontier}
A useful next model should handle a family of interventions, not just reproduce an unchanged probability table. For source data \(\sigma\), a source instrument \(\mathfrak I_{a|x}\), and a typed reduction \(\Pi\), the desired relation is schematically
\begin{equation}
\Pi\!\left(\mathfrak I_{a|x}(\sigma)\right)
 =\mathcal I_{a|x}\!\left(\Pi(\sigma)\right).
\end{equation}
Domains, positivity, normalization, composition, and the meaning of each outcome must be supplied. The equation is an experimental completion contract, not a universal theorem established here.

The relevant family includes phase changes, imperfect one-arm monitoring, absorption, null records, ordinary pointer readout, eraser conditioning, and coherent reversal. Entangled-pair tests add the requirement that allowed nonselective remote operations preserve the correct local marginals. Temporal tests add retained environmental memory and causal ordering.

Apparatus settings such as slit geometry, efficiency, timing, and preparation are legitimate experimental inputs. They are not fundamental constants that a theory must predict without being told which experiment was built. The scientific demand is one law for the response to those controls, not a separate fit to each desired output.

If objective actualization is claimed, an event law must accompany the operational instruments. The canonical MTT recorder and exact non-entailment result clarify the distinction. They are not discarded because the larger physical comparison remains to be done.

\section{Conclusion}
The interpretation developed here is one world of locally occurring events constrained by coherent source relations. Superposition is phase-sensitive plurality, not necessarily multiple actual worlds. Entanglement is a joint-state property, not mere common ancestry. Measurement is an ordinary physical process producing records. Locality requires a qualified profile. The shared circle, fixed points, and source geometry are proposed structures of comparison and admissibility, not automatic channels for instantaneous influence.

The account makes Wigner's friend intelligible without placing Wigner outside reality: a retained record can be communicated, while coherent reversal is a different experiment. The claim that actuality need not automatically destroy all coherence is stated as a commitment requiring a physical completion, not hidden inside the word projection.

The exact MTT recorder, free local net, descent, leakage, and ontology non-entailment results give this interpretation mathematical support at their declared scopes. The remaining achievement would be to bind the richer apparatus family and, if demanded, one actual-event law to the same selected source. Conceptual clarity is useful now; empirical and ontological completion remain distinguishable accomplishments.

% BEGIN MTT MANAGED COMPUTATIONAL EVIDENCE
\section*{Computational Evidence and Reproducibility}
\addcontentsline{toc}{section}{Computational Evidence and Reproducibility}
\begingroup
\raggedright
The curated source is \url{https://github.com/PeterNero/mtt-results-repro},
inspected at revision \texttt{db79bcebec12da558eef22f11ae7bbe8171fb595}.
Each identifier below resolves under
\texttt{release/results/<identifier>/artifact.json}.
The accompanying \texttt{evidence-manifest.json} records hashes, source roles,
and the status snapshot used for this manuscript.

\begin{itemize}
\item \texttt{q79\_fock\_output\_measure} and \texttt{q79\_projective\_record\_descent}: canonical first-record law, instrument descent, and the exact checkpoint.
\item \texttt{q79\_minimal\_recorder\_action} and \texttt{q79\_continuum\_recorder\_compiler}: recorder construction and its scoped extension contract.
\item \texttt{q79\_free\_dirac\_car\_net}: the selected free-field locality result.
\item \texttt{local\_cech\_hilbert\_descent} and \texttt{local\_closure\_repair\_descent}: compatible local structures, not automatic physical actualization.
\item \texttt{causal\_base\_constraint\_compression\_leakage\_packet}: the exact leakage identity and its stated locality boundary.
\item \texttt{q79\_ontology\_nonentailment}: the checkpoint law in Eq.~\eqref{eq:448} and its ontology scope.
\end{itemize}

The canonical first-record status is additionally controlled by the current
\texttt{ENC.QM.BORN} and \texttt{ENC.QM.MEAS} records; older artifact status
strings are not used to reopen their closed canonical domain.
The free-net statement is controlled by \texttt{ENC.QFT.AQFT}.
These records were retrieved on 6 September 2026 from the model generated on
2 September 2026. The model hash is recorded in the manifest.

Thirty-one small supplementary checks reproduce the explanatory examples.
They are not new selected-source calculations or experimental data. From
this paper directory, run:
\begin{quote}
\small
\begin{verbatim}
python -B examples/check_explanatory_examples_v2.py
python -B examples/check_ontology_examples_v1.py
\end{verbatim}
\end{quote}
Both use only Python's standard library.
Rational examples use exact fractions; phase, rotation, and limiting examples
include floating-point checks and must not be mistaken for interval-certified
continuum proofs.

No calculation-capsule DOI is invented. A later release can add one after the
supplement and cited repository snapshot have been curated for deposition.
The source and discussion are independently readable without a running Kernel.
\endgroup
% END MTT MANAGED COMPUTATIONAL EVIDENCE

\appendix
\section{What the supplementary checks cover}
The first eighteen checks cover monitor completeness; preservation of path
populations; a right-only input; balanced output probabilities; imperfect null
records; absorption accounting; eraser averaging; a complex marker-overlap
convention; independent-system conditioning; entangled remote marginals;
forward disturbance; earlier-marginal invariance; temporal bounds and changed
protocols; independent count moments; number-state correlations; an ideal
Zeno sequence; and the classical CHSH assignments.

The next thirteen cover reversible memory coupling; identical pointer
marginals of coherent and mixed joint states; their different coherence
readouts; incompatible readout operators; reversal with and without external
copies; partial record overlap; the 448-atom partition; basis-dependent
coefficient counting; ensemble nonuniqueness; and a conditional Poisson race.

The models intentionally expose assumptions. They do not reconstruct every
systematic error of the cited experiments, select a physical HYM endpoint,
derive a universal apparatus action, or validate a new interpretation solely
because a program exits successfully.

\section{Relationship to the earlier interpretation}
The earlier causal-base essay asks how closure repair, geometry, effective
fields, and admissibility could form one philosophical account. This paper
asks how that account should speak about concrete quantum experiments. The
companion is not superseded. The present text adds a qualified locality
signature, one-arm and temporal protocol distinctions, particle-number
clarifications, and an explicit local-first Wigner account.

Expressions such as global admissibility are used here for compatibility
conditions on descriptions where those conditions are meaningful. They are
not a declaration that a completed future or an ontically fundamental
universal state exists. This clarification preserves the process reading
while leaving the relevant extension and dynamics problems explicit.

\begin{thebibliography}{99}
\raggedright
\bibitem{Bach}
R. Bach, D. Pope, S.-H. Liou, and H. Batelaan,
\emph{Controlled double-slit electron diffraction},
New Journal of Physics 15, 033018 (2013).
\url{https://arxiv.org/abs/1210.6243}.
\bibitem{Grangier}
P. Grangier, G. Roger, and A. Aspect,
\emph{Experimental evidence for a photon anticorrelation effect on a beam
splitter: a new light on single-photon interferences},
Europhysics Letters 1, 173--179 (1986).
\url{https://doi.org/10.1209/0295-5075/1/4/004}.
\bibitem{Glauber}
R. J. Glauber, \emph{The quantum theory of optical coherence},
Physical Review 130, 2529--2539 (1963).
\url{https://doi.org/10.1103/PhysRev.130.2529}.
\bibitem{Englert}
B.-G. Englert, \emph{Fringe visibility and which-way information: an inequality},
Physical Review Letters 77, 2154--2157 (1996).
\url{https://doi.org/10.1103/PhysRevLett.77.2154}.
\bibitem{ElitzurVaidman}
A. C. Elitzur and L. Vaidman, \emph{Quantum mechanical interaction-free
measurements}, Foundations of Physics 23, 987--997 (1993).
\url{https://arxiv.org/abs/hep-th/9305002}.
\bibitem{HongOuMandel}
C. K. Hong, Z. Y. Ou, and L. Mandel,
\emph{Measurement of subpicosecond time intervals between two photons by
interference}, Physical Review Letters 59, 2044--2046 (1987).
\url{https://doi.org/10.1103/PhysRevLett.59.2044}.
\bibitem{Werner}
R. F. Werner, \emph{Quantum states with Einstein-Podolsky-Rosen correlations
admitting a hidden-variable model}, Physical Review A 40, 4277--4281 (1989).
\url{https://doi.org/10.1103/PhysRevA.40.4277}.
\bibitem{Bell}
J. S. Bell, \emph{On the Einstein Podolsky Rosen paradox},
Physics Physique Fizika 1, 195--200 (1964).
\url{https://doi.org/10.1103/PhysicsPhysiqueFizika.1.195}.
\bibitem{Hensen}
B. Hensen et al., \emph{Loophole-free Bell inequality violation using
electron spins separated by 1.3 kilometres}, Nature 526, 682--686 (2015).
\url{https://arxiv.org/abs/1508.05949}.
\bibitem{Jacques}
V. Jacques et al., \emph{Experimental realization of Wheeler's delayed-choice
Gedanken experiment}, Science 315, 966--968 (2007).
\url{https://arxiv.org/abs/quant-ph/0610241}.
\bibitem{Kim}
Y.-H. Kim, R. Yu, S. P. Kulik, Y. Shih, and M. O. Scully,
\emph{Delayed choice quantum eraser},
Physical Review Letters 84, 1--5 (2000).
\url{https://arxiv.org/abs/quant-ph/9903047}.
\bibitem{Pollock}
F. A. Pollock, C. Rodr\'iguez-Rosario, T. Frauenheim, M. Paternostro,
and K. Modi, \emph{Operational Markov condition for quantum processes},
Physical Review Letters 120, 040405 (2018).
\url{https://arxiv.org/abs/1801.09811}.
\bibitem{LeggettGarg}
A. J. Leggett and A. Garg, \emph{Quantum mechanics versus macroscopic
realism: is the flux there when nobody looks?},
Physical Review Letters 54, 857--860 (1985).
\url{https://doi.org/10.1103/PhysRevLett.54.857}.
\bibitem{Robertson}
H. P. Robertson, \emph{The uncertainty principle},
Physical Review 34, 163--164 (1929).
\url{https://doi.org/10.1103/PhysRev.34.163}.
\bibitem{AharonovBohm}
Y. Aharonov and D. Bohm, \emph{Significance of electromagnetic potentials
in the quantum theory}, Physical Review 115, 485--491 (1959).
\url{https://doi.org/10.1103/PhysRev.115.485}.
\bibitem{FewsterVerch}
C. J. Fewster and R. Verch, \emph{Quantum fields and local measurements},
Communications in Mathematical Physics 378, 851--889 (2020).
\url{https://arxiv.org/abs/1810.06512}.
\bibitem{AbramskyBrandenburger}
S. Abramsky and A. Brandenburger,
\emph{The sheaf-theoretic structure of non-locality and contextuality},
New Journal of Physics 13, 113036 (2011).
\url{https://arxiv.org/abs/1102.0264}.
\bibitem{FrauchigerRenner}
D. Frauchiger and R. Renner, \emph{Quantum theory cannot consistently
describe the use of itself}, Nature Communications 9, 3711 (2018).
\url{https://arxiv.org/abs/1604.07422}.
\bibitem{Bong}
K.-W. Bong et al., \emph{A strong no-go theorem on the Wigner's friend
paradox}, Nature Physics 16, 1199--1205 (2020).
\url{https://arxiv.org/abs/1907.05607}.
\bibitem{NeroConstraint}
P. Nero, \emph{A philosophical interpretation of causal-base constraint
realism}, working manuscript v2 (2026), companion interpretive source.
\url{https://github.com/PeterNero/mtt-causal-base-constraint-fiber/blob/main/PHILOSOPHICAL_INTERPRETATION_OF_CONSTRAINT_REALISM_v2.md}.
\bibitem{NeroFoundation}
P. Nero, \emph{Modal Triplet Theory: Foundations}, v9 (2026).
\url{https://doi.org/10.5281/zenodo.21655367}.
\bibitem{NeroProto}
P. Nero, \emph{The Proto-Spinor: Conditional Spinorial Closure and the
q79 Interface}, v7 (2026).
\url{https://doi.org/10.5281/zenodo.21655396}.
\bibitem{NeroBell}
P. Nero, \emph{Modal Fixed Points, Bell's Beables, and the Limits of
Factorization: Upper-Local Dynamics, Nonseparable States, and Operational
No-Signaling in MTT}, v2 (2026).
\url{https://doi.org/10.5281/zenodo.21665825}.
\bibitem{NeroTemporal}
P. Nero, \emph{Temporal Bell Inequalities and Global Consistency in Modal
Triplet Theory: A Conditional Sequential-Instrument Analysis}, v3 (2026).
\url{https://doi.org/10.5281/zenodo.21657161}.
\bibitem{NeroMeasurement}
P. Nero, \emph{Why Decoherence Cannot Replace Measurement:
Outcome-Resolved Dynamics Without a Fundamental Measurement Postulate
in Modal Triplet Theory}, v2 (2026).
\url{https://doi.org/10.5281/zenodo.21666019}.
\bibitem{NeroQM}
P. Nero, \emph{MTT quantum-mechanics source proofs},
canonical recorder and conditional actualization sources (2026).
\url{https://github.com/PeterNero/mtt-qm-source-proof}.
\bibitem{NeroRepair}
P. Nero, \emph{Cohesive Closure Repair and Its Hodge, Kernel, and Projection
Shadows: From Nonlinear Defects to Tangent Semigroups, with the
Physical-Action Boundary}, unpublished canonical source project (2026).
Project identifier:
\texttt{cohesive-closure-repair-and-its-hodge-kernel-and-projection-shadows}.
The underlying structural artifacts are separately available in \cite{NeroResults}.
\bibitem{NeroResults}
P. Nero, \emph{MTT results reproducibility repository},
revision db79bcebec12da558eef22f11ae7bbe8171fb595 (2026).
\url{https://github.com/PeterNero/mtt-results-repro/tree/db79bcebec12da558eef22f11ae7bbe8171fb595/release/results}.
\end{thebibliography}
\end{document}
